\documentclass[12pt]{article}
\usepackage[utf8]{inputenc}
\usepackage[T1]{fontenc}
\usepackage[margin=1in]{geometry}
\usepackage{setspace}
\usepackage{xcolor}
\usepackage[colorlinks=true, urlcolor=blue!60!black, linkcolor=blue!60!black]{hyperref}
\setlength{\parindent}{2em}
\setlength{\parskip}{1.25em}
\renewcommand{\baselinestretch}{1.0}
\begin{document}
\begin{center}
{\LARGE \textit{Learning Resources, Inc. v. Trump}, 607 U.S. \_\_\_ (2026)\par}
\vspace{0.8em}
{\large Case Brief --- Four-Page Edition\par}
\end{center}
\vspace{1em}

\section{Memory Jogger}

The Court held 6--3 that the International Emergency Economic Powers Act's grant of authority to ``regulate \ldots{} importation'' does not authorize the President to impose tariffs, while fracturing sharply over the Major Questions Doctrine and the Nondelegation Doctrine (\href{https://www.supremecourt.gov/opinions/25pdf/24-1287_4gcj.pdf}{Slip Op.}).

\section{Detailed Case Facts}

In early 2025 President Trump declared national emergencies and invoked IEEPA, which lets the President ``regulate'' or ``prohibit'' importation to deal with an ``unusual and extraordinary threat'' from abroad. To address fentanyl and other drugs, he imposed ``drug trafficking'' tariffs of 25\% on most Canadian and Mexican imports and 10\% (later 20\%) on most Chinese imports. Citing ``large and persistent'' trade deficits, he then imposed worldwide ``reciprocal'' tariffs of at least 10\%, escalating the total rate on most Chinese goods to 145\%. The government collected more than \$200 billion.

Learning Resources, Inc. and hand2mind, Inc., family-owned makers of children's educational products manufactured largely abroad, faced duties rising from \$2.3 million in 2024 to over \$100 million projected for 2025. In \textit{V.O.S. Selections}, five small businesses and twelve States also challenged both sets of tariffs.

\section{Procedural History}

The District of Columbia district court refused to transfer the case to the Court of International Trade (CIT) and preliminarily enjoined the tariffs. Separately, the CIT granted summary judgment to the \textit{V.O.S.} plaintiffs, and the en banc Federal Circuit affirmed, holding that IEEPA did not authorize tariffs ``unbounded in scope, amount, and duration.'' The Supreme Court consolidated the cases, heard argument on November 5, 2025, and decided them on February 20, 2026.

\section{Judicial Votes}

\textbf{6--3} against the tariffs, with seven opinions.
\begin{itemize}
\item \textbf{Opinion of the Court (Parts I, II--A--1, II--B):} Chief Justice Roberts, joined by Justices Sotomayor, Kagan, Gorsuch, Barrett, and Jackson. Parts II--A--2 and III (Major Questions Doctrine) drew only Gorsuch and Barrett.
\item \textbf{Concurrences:} Justices Gorsuch and Barrett (each in full); Justice Kagan, joined by Sotomayor and Jackson, and Justice Jackson (each in part and in the judgment).
\item \textbf{Dissents:} Justice Kavanaugh, joined by Justices Thomas and Alito; Justice Thomas separately.
\end{itemize}

\section{Holding}

IEEPA does not authorize the President to impose tariffs. Article I vests the taxing and tariff power in Congress, the President has no inherent peacetime tariff power, and ``regulate \ldots{} importation'' does not encompass the distinct power to levy revenue-raising duties. The Federal Circuit was affirmed; the \textit{Learning Resources} judgment was vacated for lack of jurisdiction because the CIT's jurisdiction is exclusive. The Major Questions Doctrine analysis is plurality only, and the Court left open non-tariff IEEPA powers, Title 19 tariffs, reviewability of emergency declarations, and refunds.

\section{Analysis of Opinions}

\textbf{Roberts (Court and plurality).} Tariffs are ``a branch of the taxing power,'' and the government relied solely on IEEPA. For the plurality, the Major Questions Doctrine applied: Congress delegates tariff power expressly with caps and limits, no President had used IEEPA for tariffs in nearly 50 years, and the economic stakes were vast. For the majority, IEEPA never mentions tariffs, ``regulate'' does not ordinarily include taxation, reading it so would make export regulation unconstitutional under the Export Clause, and the neighboring verbs describe sanctions, not taxes.

\textbf{Gorsuch (concurring).} The Major Questions Doctrine is a Constitution-enforcing clear-statement rule rooted in Article I; because Congress can retrieve delegated power after \textit{Chadha} only over a veto, it is ``pro-Congress.''

\textbf{Barrett (concurring).} The doctrine is ordinary contextual textualism, not a ``strong-form substantive canon'' imposing a ``clarity tax.''

\textbf{Kagan (concurring in part).} Ordinary construction suffices: of IEEPA's 99 authorized actions, none of the other 98 raises revenue, so tariffs would be the ``odd man out.''

\textbf{Jackson (concurring in part).} The 1941 and 1977 legislative reports show Congress meant only to let the President ``control or freeze property transactions where a foreign interest is involved.''

\textbf{Kavanaugh (dissenting).} If quotas and embargoes regulate importation, ``how are tariffs not a means to regulate importation?'' He relied on \textit{Yoshida} and \textit{Algonquin}, rejected a ``magic-words test,'' would exempt foreign affairs from the doctrine, and warned that refunds will be a ``mess.''

\textbf{Thomas (dissenting).} Nondelegation limits only ``core legislative power'' over private rights; ``[i]mporting is a matter of privilege.''

\section{Examples: Future Applications}

\begin{enumerate}
\item \textbf{Same-Side (Invalid): Carbon Border Tax.} A President declares a climate emergency and uses IEEPA to levy a 30\% tax on imported cars and steel. Like these tariffs, it uses ``regulate'' to raise revenue without express authorization.
\item \textbf{Same-Side (Invalid): SEC Transactions Tax.} The SEC uses its power to ``regulate'' securities trading to impose a 1\% tax on trades. The power to regulate does not include the power to tax.
\item \textbf{Opposite-Side (Valid): Asset Freezes and Embargoes.} The President freezes an adversary's assets and bans its telecommunications imports. IEEPA expressly authorizes blocking and prohibiting, without taxation.
\item \textbf{Opposite-Side (Valid): Title 19 Tariffs.} A 15\% surcharge for 120 days under Section 122 of the Trade Act, or steel duties after a Section 232 investigation, rests on express delegations with caps and procedures.
\item \textbf{Fence-Sitter: Licensing Fees or Wartime Tariffs.} The President imposes a cost-recovery ``regulatory licensing fee'' on imported AI models, or tariffs during an authorized armed conflict. Challengers see a disguised tariff; the Executive invokes the difference between fees and taxes and overlapping Article II powers.
\end{enumerate}

\section{Critique}

\textbf{Textualist and originalist views.} Abbe Gluck described an open ``canon war'' between Gorsuch's external clear-statement rule and Barrett's fear of an atextual ``clarity tax.'' Thomas's attempt to confine nondelegation to deprivations of life, liberty, or property drew criticism, including from Gorsuch, as inviting executive aggrandizement.

\textbf{Progressive and administrative-law views.} Ronald Levin argued that the decision left the doctrine ``more tenuous,'' since only three Justices joined it and Kagan showed ordinary interpretation sufficed. Michael Dreeben noted unresolved questions about wartime tariffs, reviewability of emergency declarations, and Congress's willingness to reclaim powers after \textit{Chadha}.

\section{Key Quotations}

\begin{quote}
``Based on two words separated by 16 others in Section 1702(a)(1)(B) of IEEPA---`regulate' and `importation'---the President asserts the independent power to impose tariffs on imports from any country, of any product, at any rate, for any amount of time. Those words cannot bear such weight.''
\end{quote}
\noindent --- Chief Justice Roberts, Opinion of the Court, slip op. at 5.

\begin{quote}
``There is no major questions exception to the major questions doctrine.''
\end{quote}
\noindent --- Chief Justice Roberts, Plurality Opinion, slip op. at 13.

\begin{quote}
``For all those reasons, straight-up statutory construction resolves this case for me; I need no major-questions thumb on the interpretive scales.''
\end{quote}
\noindent --- Justice Kagan, concurring in part and in the judgment, slip op. at 7.

\end{document}
